\section{Introduction}
\label{sec:introduction}

When one event is described immediately after another, comprehenders often infer more than temporal succession: the first event may be understood as explaining why the second occurred. This graded causal implication supports ordinary narrative understanding, explanation, and commonsense reasoning. A language model that is sensitive to the distinction should separate genuinely explanatory factors from details that are merely compatible with an event. Yet a model's prompted answer need not transparently reveal whether it represents that distinction. Answer-token priors, scale presentation, and late-stage readout can all distort an otherwise informative internal signal.

This representation--behavior distinction parallels recent work on event plausibility. Lepori et al. show that linear directions in language-model activations can distinguish modal categories and align with fine-grained human judgments even when output probabilities provide an unreliable behavioral measure \citep{lepori2026fantasy}. Causal implication presents a related but more relational problem: the model must evaluate whether a stated factor explains a particular event, rather than assess the plausibility of one sentence in isolation. Temporal proximity is a particularly important alternative explanation because causes normally precede their effects, but temporal order alone does not establish an explanatory relation.

Before asking how causal implication is represented, we first test whether it is visible in unconstrained next-token behavior. Following the ignorance-continuation logic of \citet{pan-kehler-2025-pragmatic}, an inferred explanation should make a subsequent claim that the reason is unknown less expected. A complementary diagnostic compares the probability of the same factor after \emph{because} versus after the explicitly incidental continuation \emph{and it so happened that}. Together, these tests avoid requiring the model to interpret a metalinguistic Yes/No question or numeric scale.

We therefore analyze causal implication at four levels. First, we compare implicit eliciture signatures from the ignorance and explicit-happenstance diagnostics with direct probabilistic Yes/No and Likert judgments. Second, we compare prompted behavior with linear probes and paired mean-difference directions, and test whether causal and temporal representations transfer across domains. Third, we ask whether the resulting continuous representation scores track probabilistic Yes/No judgments, probabilistic Likert judgments, and an external graded rating. Fourth, we distinguish two probe-relative decision failures: \emph{suppression}, in which the causal probe decodes a factor as explanatory but the model rejects it, and \emph{post-hoc attribution}, in which the probe decodes a factor as non-explanatory but the model accepts it. We test whether these failures are predictable from pre-answer activations and whether subtracting development-selected failure directions changes the final decision.

Our experiments are organized around the following research questions:

\begin{samepage}
\noindent\textbf{RQ1: Behavioral eliciture diagnostics.} Do causal critical items, relative to matched controls, increase the surprisal of a subsequent ignorance claim and increase the probability of the explanatory factor after \emph{because} relative to \emph{and it so happened that}? How consistently does this implicit sensitivity appear in direct probabilistic Yes/No and Likert judgments?
\end{samepage}

\begin{samepage}
\noindent\textbf{RQ2: Decodability and transfer.} Is graded causal implication linearly decodable across pretrained model families, and is the decoded causal signal distinct from a direction learned for temporal relations? How does internal decodability compare with prompted Yes/No and Likert behavior?
\end{samepage}

\begin{samepage}
\noindent\textbf{RQ3: Correlation with graded judgments.} Do raw causal probe logits and CAA projections track item-level graded judgments from the target model and an external rater, both across individual items and in critical-minus-control pair differences?
\end{samepage}

\Needspace{6\baselineskip}
\begin{samepage}
\noindent\textbf{RQ4: Predicting and reducing failure modes.} Can probe-relative suppression and post-hoc attribution be predicted from pre-answer activations and reduced by format-balanced, development-selected residual-stream steering without substantially damaging probe--output agreements?
\end{samepage}
